The Python implementation is reproduced from the repository root (the
directory containing \path{main.tex} and \path{common.py}) as follows.
Each task has an entry-point script, a report fragment, implementation notes,
and an \path{outputs/} directory.

\begin{enumerate}[leftmargin=*,itemsep=0.5em]
\item \textbf{Check the supplied data.}
Place \path{Test_Image_2026.bmp} in \path{dataset_2026/}, alongside the seven
class directories \texttt{0, 4, 7, 8, A, D, H}.  Each class contains 1,016
PNG images (7,112 images in total), with source dimensions $128\times128$.
Keep the BMP and its derived crops outside the labelled directories: they are
reserved for Figure~1 inference.

\item \textbf{Create the software environment.}
Install Conda, then execute:
{\small
\begin{verbatim}
conda env create -f task_7/environment.yml
source "$(conda info --base)/etc/profile.d/conda.sh"
conda activate cvtask
\end{verbatim}
}
If \texttt{cvtask} already exists, replace the creation command with
\texttt{conda env update -n cvtask -f task\_7/environment.yml}.
The environment file pins Python~3.8, PyTorch~2.0.0, torchvision~0.15.1~\cite{torchvision},
CUDA runtime~11.8, NumPy~1.24.2, Pillow~9.4.0, and Matplotlib~3.7.1.
The measured runtime was Ubuntu~20.04, Python~3.8.10, and
PyTorch~2.0.0+cu118 on an NVIDIA GeForce RTX~4080 SUPER.
Tasks~1--6 run on the CPU.  Task~7 automatically selects CUDA if available,
otherwise the CPU; prefixing its command with
\texttt{CUDA\_VISIBLE\_DEVICES=""} forces CPU execution.  Timings are
device-dependent.  Plots use the Agg backend and require no display server.

\item \textbf{Run every image-processing stage in order.}
{\small
\begin{verbatim}
python task_1/task_1.py
python task_2/task_2.py
python task_3/task_3.py
python task_4/task_4.py
python task_5/task_5.py
python task_6/task_6.py
\end{verbatim}
}
These scripts create or replace files in their fixed
\path{task_N/outputs/} directories.  Task~1 reads the BMP and saves
\path{task1_percentile_stretch_p02_p98.png} for Tasks~2 and~3.
Task~2 is the filter-size comparison branch; Task~3 uses the Task~1 enhanced
image and original BMP to save \path{task_3/outputs/task3_subimage.png}.
Task~4 produces \path{task_4/outputs/task4_binary.png} and the Otsu comparison.
Tasks~5 and~6 read the Task~3 grayscale image and Task~4 binary results.
Task~5 writes outline figures and JSON contour records.
Task~6 writes \path{task6_components.csv} and ten pairs of
\path{task_6/outputs/chars/char_*_gray.png} and \path{char_*_bin.png} crops;
these must exist before Task~7 to obtain the Figure~1 recognition results.

\item \textbf{Run the complete learning experiment.}
Use a separate directory to preserve the report's supplied results:
{\small
\begin{verbatim}
python task_7/task_7.py --data dataset_2026 \
  --output task_7/outputs_reproduction
\end{verbatim}
}
The default seed is 20261013, batch size 128, tuning limit 18 epochs,
early-stopping patience 5, and data-loader workers 0.  The full run performs
the $32/48$-pixel and $10^{-3}/3\times10^{-4}$ learning-rate grid, chooses
models on validation data, refits on the full 75\% training pool using the
selected validation epoch, then evaluates the 25\% test partition and the
target crops.  It uses one CNN and three MLP regression members.
For a short pipeline check, use:
{\small
\begin{verbatim}
python task_7/task_7.py --data dataset_2026 \
  --smoke --output task_7/outputs_smoke
\end{verbatim}
}
Smoke mode uses one input size, one learning rate, at most two tuning epochs,
and one MLP; it is not the reported experiment.  Always specify the smoke
directory because the default output is \path{task_7/outputs/}.
The main artifacts are \path{split_summary.json} and
\path{split_manifest.csv} (partition audit), \path{tuning_results.csv} and
\path{tuning_curves.png} (model selection), \path{task7_metrics.json} and
\path{method_comparison.csv} (results and timing), two confusion-matrix PNGs,
\path{*.pt} checkpoints, and \path{target_predictions.json}.
The report reads its formal figures from \path{task_7/outputs/}; independent
reproduction and smoke runs leave these unchanged.  Replacing the formal
directory requires updating the report's numerical tables to match the new
metrics.  Exact numerical agreement and timing across different software or
hardware are not guaranteed.

To repeat target inference from the supplied models without retraining, copy
\path{task_7/outputs/} to a new \path{task_7/outputs_inference_check/}
directory, then run:
{\small
\begin{verbatim}
python task_7/task_7.py --data dataset_2026 --inference-only \
  --output task_7/outputs_inference_check
\end{verbatim}
}
This mode needs both crop variants, the dataset, split manifest, and copied
checkpoints.  It fits foreground geometry from training rows only and updates
target inference while preserving test metrics and the grayscale baseline.

\item \textbf{Compile the report.}
Install XeLaTeX, latexmk, and the recommended/extra LaTeX packages.
On Ubuntu, run:
{\small
\begin{verbatim}
sudo apt-get install latexmk texlive-xetex texlive-latex-extra \
  texlive-latex-recommended texlive-fonts-recommended \
  curl cabextract fontconfig
\end{verbatim}
}
Omit \texttt{sudo} when running as root.  The provided font installer verifies
and installs the actual Arial family (regular, bold, italic, and bold italic)
under \path{~/.local/share/fonts/msttcorefonts/} by default, without root
privileges; \texttt{REPORT\_FONT\_DIR} can override the directory.
With all task figures present, execute:
{\small
\begin{verbatim}
bash scripts/install_report_fonts.sh
fc-match Arial
latexmk -xelatex -interaction=nonstopmode -halt-on-error \
  -file-line-error main.tex
\end{verbatim}
}
The font check must resolve to Arial.  The build also runs BibTeX for the
references.  The output is \path{main.pdf}.  The project \path{.latexmkrc} removes auxiliary
build files after a successful build while retaining the PDF.
\end{enumerate}

\noindent\path{README.md} contains the full directory dependency table,
artifact inventory, version-check command, and build troubleshooting command.
